\documentclass{article}
\usepackage{fancyhdr}
\usepackage[table]{xcolor}
\usepackage[margin=1in]{geometry}
\usepackage{caption}
\usepackage{hyperref}
\usepackage{graphicx}
\usepackage{float}
\usepackage{enumitem}
\usepackage{array}
\usepackage{tikz}

\geometry{a4paper, margin=1in}

\title{Water Transfer System Improvement Report}
\author{Your Name}
\date{\today}

\begin{document}


\section{Summary of the 16 Concepts}

This section provides a brief overview of the sixteen concepts developed for the efficient water transfer system. Each concept integrates various combinations of pumps, energy sources, and storage solutions to enhance water delivery from the elevated tank to the garden area.

\begin{enumerate}[label=\textbf{\arabic*.}]
    \item \textbf{Modified System Using a 2-Inch Pipe}: Utilizes a wider 2-inch pipe to reduce friction losses, significantly increasing flow rate and pressure without additional energy sources.
    
    \item \textbf{Gravity System by Raising the Tank}: Elevates the primary tank to increase gravitational pressure, enhancing natural water flow and reducing the need for mechanical pumps.
    
    \item \textbf{Manual Pump System}: Incorporates a manual pump to boost water pressure and flow, relying on human effort to operate the pump.
    
    \item \textbf{Pump System with Raised Tank}: Combines a manual pump with a raised secondary tank to leverage both manual pumping and gravitational pressure for improved water delivery.
    
    \item \textbf{Pump System with Pressure Tank}: Integrates a manual pump with a pressure tank to maintain consistent water flow and pressure, reducing the need for continuous pumping.
    
    \item \textbf{Solar Pump System}: Employs a solar-powered pump to automate water transfer, minimizing manual intervention and utilizing renewable energy.
    
    \item \textbf{Solar-Powered Pump with Raised Tank System}: Combines a solar-powered pump with an elevated secondary tank to enhance water pressure through both solar energy and gravity.
    
    \item \textbf{Solar-Powered Pump with Pressure Tank System}: Integrates a solar-powered pump with a pressure tank to ensure stable and controlled water flow using renewable energy sources.
    
    \item \textbf{Battery-Powered Pump System}: Uses a battery-operated pump to provide consistent water transfer regardless of sunlight availability, offering more reliable performance than solar-only systems.
    
    \item \textbf{Battery-Powered Pump with Pressure Tank System}: Combines a battery-powered pump with a pressure tank to deliver steady water flow and pressure, enhancing reliability and reducing manual effort.
    
    \item \textbf{Battery-Powered Pump with Raised Tank System}: Utilizes a battery-powered pump to transfer water to a raised tank, leveraging battery power and gravity for efficient water delivery.
    
    \item \textbf{Solar + Battery + Pressure Tank System}: Integrates solar power, battery storage, and a pressure tank to create a hybrid system that ensures continuous and reliable water flow using multiple energy sources.
    
    \item \textbf{Solar + Battery-Powered Pump with Raised Tank System}: Combines solar energy and battery power to operate a pump that fills a raised tank, utilizing both renewable energy and gravitational force for steady water flow.
    
    \item \textbf{Grid-Powered Pump System}: Employs a pump powered by the electrical grid to achieve the highest flow rate and pressure, ensuring robust performance but relying on non-renewable energy sources.
    
    \item \textbf{Grid-Powered Pump with Pressure Tank System}: Integrates a grid-powered pump with a pressure tank to maintain consistent water flow and pressure, optimizing performance through reliable grid energy.
    
    \item \textbf{Grid-Powered Pump with Raised Tank System}: Uses a grid-powered pump to transfer water to an elevated tank, leveraging both electrical power and gravity to ensure efficient and steady water delivery.
\end{enumerate}

\noindent Each concept addresses the core functions of activation, flow rate adjustment, water storage, transfer, flow control, and delivery, employing different combinations of energy sources and storage mechanisms to meet the project's efficiency and sustainability goals.


\end{document}
